\documentclass[11pt]{article}
% =============================================================================
%  Shared article preamble — tutorials, homework, solutions.
%  Same fonts, palette and notation as the Beamer decks (see preamble.tex).
%      \documentclass[11pt]{article}
%      % =============================================================================
%  Shared article preamble — tutorials, homework, solutions.
%  Same fonts, palette and notation as the Beamer decks (see preamble.tex).
%      \documentclass[11pt]{article}
%      % =============================================================================
%  Shared article preamble — tutorials, homework, solutions.
%  Same fonts, palette and notation as the Beamer decks (see preamble.tex).
%      \documentclass[11pt]{article}
%      \input{../../../style/preamble_article}
% =============================================================================
\usepackage[margin=1in]{geometry}
\usepackage{helvet}
\usepackage[default]{lato}
\usepackage{mathpazo}
\usepackage[utf8]{inputenc}
\usepackage{amsmath,amssymb,amsthm}
\usepackage{mathtools}
\usepackage{graphicx}
\usepackage{booktabs}
\usepackage{multirow}
\usepackage{enumitem}
\usepackage{xcolor}
\usepackage{tcolorbox}
\tcbuselibrary{skins,breakable}
\usepackage{listings}
\usepackage[ruled,vlined]{algorithm2e}
\usepackage{hyperref}
\usepackage{fancyhdr}
\usepackage{titlesec}

\definecolor{blue}{RGB}{0,114,178}
\definecolor{red}{RGB}{213,94,0}
\definecolor{yellow}{RGB}{240,228,66}
\definecolor{green}{RGB}{0,158,115}
\hypersetup{colorlinks=true, linkcolor=blue, urlcolor=blue, citecolor=blue}

\titleformat{\section}{\Large\bfseries\color{blue}}{\thesection}{0.6em}{}
\titleformat{\subsection}{\large\bfseries}{\thesubsection}{0.6em}{}
\pagestyle{fancy}\fancyhf{}
\rhead{\footnotesize ENEI -- INEI \,\textperiodcentered\, PEU-CD 2026}
\lfoot{\footnotesize Machine Learning I \& II \,\textperiodcentered\, Alexander Quispe}
\rfoot{\footnotesize\thepage}
\renewcommand{\headrulewidth}{0.3pt}

\newtcolorbox{derivation}[1]{enhanced, breakable, colback=blue!3, colframe=blue, fonttitle=\bfseries,
  title={#1}, boxrule=0.6pt, left=6pt, right=6pt, top=4pt, bottom=4pt, arc=2pt}
\newtcolorbox{claim}[1]{enhanced, breakable, colback=white, colframe=black, fonttitle=\bfseries,
  title={#1}, boxrule=0.6pt, left=6pt, right=6pt, top=4pt, bottom=4pt, arc=2pt}
\newtcolorbox{keypoint}{enhanced, breakable, colback=green!6, colframe=green, boxrule=0.6pt,
  left=6pt, right=6pt, top=4pt, bottom=4pt, arc=2pt}
\newtcolorbox{caution}{enhanced, breakable, colback=red!5, colframe=red, boxrule=0.6pt,
  left=6pt, right=6pt, top=4pt, bottom=4pt, arc=2pt}
\newtcolorbox{task}[1]{enhanced, breakable, colback=yellow!12, colframe=yellow!60!black, fonttitle=\bfseries,
  title={#1}, boxrule=0.6pt, left=6pt, right=6pt, top=4pt, bottom=4pt, arc=2pt}

\lstset{
  basicstyle=\ttfamily\small, keywordstyle=\color{blue}\bfseries,
  commentstyle=\color{green}, stringstyle=\color{red},
  showstringspaces=false, breaklines=true, frame=single, rulecolor=\color{black!30},
  numbers=none, xleftmargin=6pt, framexleftmargin=6pt, language=Python
}
\newtheorem{theorem}{Theorem}
\newtheorem{lemma}{Lemma}
\newtheorem{proposition}{Proposition}
\theoremstyle{definition}
\newtheorem{definition}{Definition}
\newtheorem{exercise}{Exercise}

\input{../../../style/macros}

% =============================================================================
\usepackage[margin=1in]{geometry}
\usepackage{helvet}
\usepackage[default]{lato}
\usepackage{mathpazo}
\usepackage[utf8]{inputenc}
\usepackage{amsmath,amssymb,amsthm}
\usepackage{mathtools}
\usepackage{graphicx}
\usepackage{booktabs}
\usepackage{multirow}
\usepackage{enumitem}
\usepackage{xcolor}
\usepackage{tcolorbox}
\tcbuselibrary{skins,breakable}
\usepackage{listings}
\usepackage[ruled,vlined]{algorithm2e}
\usepackage{hyperref}
\usepackage{fancyhdr}
\usepackage{titlesec}

\definecolor{blue}{RGB}{0,114,178}
\definecolor{red}{RGB}{213,94,0}
\definecolor{yellow}{RGB}{240,228,66}
\definecolor{green}{RGB}{0,158,115}
\hypersetup{colorlinks=true, linkcolor=blue, urlcolor=blue, citecolor=blue}

\titleformat{\section}{\Large\bfseries\color{blue}}{\thesection}{0.6em}{}
\titleformat{\subsection}{\large\bfseries}{\thesubsection}{0.6em}{}
\pagestyle{fancy}\fancyhf{}
\rhead{\footnotesize ENEI -- INEI \,\textperiodcentered\, PEU-CD 2026}
\lfoot{\footnotesize Machine Learning I \& II \,\textperiodcentered\, Alexander Quispe}
\rfoot{\footnotesize\thepage}
\renewcommand{\headrulewidth}{0.3pt}

\newtcolorbox{derivation}[1]{enhanced, breakable, colback=blue!3, colframe=blue, fonttitle=\bfseries,
  title={#1}, boxrule=0.6pt, left=6pt, right=6pt, top=4pt, bottom=4pt, arc=2pt}
\newtcolorbox{claim}[1]{enhanced, breakable, colback=white, colframe=black, fonttitle=\bfseries,
  title={#1}, boxrule=0.6pt, left=6pt, right=6pt, top=4pt, bottom=4pt, arc=2pt}
\newtcolorbox{keypoint}{enhanced, breakable, colback=green!6, colframe=green, boxrule=0.6pt,
  left=6pt, right=6pt, top=4pt, bottom=4pt, arc=2pt}
\newtcolorbox{caution}{enhanced, breakable, colback=red!5, colframe=red, boxrule=0.6pt,
  left=6pt, right=6pt, top=4pt, bottom=4pt, arc=2pt}
\newtcolorbox{task}[1]{enhanced, breakable, colback=yellow!12, colframe=yellow!60!black, fonttitle=\bfseries,
  title={#1}, boxrule=0.6pt, left=6pt, right=6pt, top=4pt, bottom=4pt, arc=2pt}

\lstset{
  basicstyle=\ttfamily\small, keywordstyle=\color{blue}\bfseries,
  commentstyle=\color{green}, stringstyle=\color{red},
  showstringspaces=false, breaklines=true, frame=single, rulecolor=\color{black!30},
  numbers=none, xleftmargin=6pt, framexleftmargin=6pt, language=Python
}
\newtheorem{theorem}{Theorem}
\newtheorem{lemma}{Lemma}
\newtheorem{proposition}{Proposition}
\theoremstyle{definition}
\newtheorem{definition}{Definition}
\newtheorem{exercise}{Exercise}

% =============================================================================
%  Notation — one convention for all lectures
%
%  scalars      x, y, w           plain italic
%  vectors      \vx \vy \vw       bold lower-case
%  matrices     \vX \vW \vSigma   bold upper-case
%  parameters   \vbeta \vtheta    bold Greek
%  transpose    \vX\T             \top
% =============================================================================
\newcommand{\R}{\mathbb{R}}
\newcommand{\E}{\mathbb{E}}
\newcommand{\Prob}{\mathbb{P}}
\DeclareMathOperator{\Var}{Var}
\DeclareMathOperator{\Cov}{Cov}
\DeclareMathOperator{\tr}{tr}
\DeclareMathOperator{\diag}{diag}
\DeclareMathOperator{\rank}{rank}
\DeclareMathOperator{\sign}{sign}
\DeclareMathOperator{\col}{col}
\DeclareMathOperator{\RSS}{RSS}
\DeclareMathOperator*{\argmin}{arg\,min}
\DeclareMathOperator*{\argmax}{arg\,max}
\newcommand{\norm}[1]{\left\lVert #1 \right\rVert}
\newcommand{\abs}[1]{\left\lvert #1 \right\rvert}
\newcommand{\inner}[2]{\left\langle #1,\, #2 \right\rangle}
\newcommand{\ind}[1]{\mathbf{1}\!\left\{#1\right\}}
\newcommand{\T}{^{\top}}
\newcommand{\inv}{^{-1}}
\newcommand{\half}{\tfrac{1}{2}}
\newcommand{\eps}{\varepsilon}
\newcommand{\Dset}{\mathcal{D}}
\newcommand{\Hset}{\mathcal{H}}
\newcommand{\Lcal}{\mathcal{L}}
\newcommand{\Ncal}{\mathcal{N}}

% bold vectors
\newcommand{\vx}{\mathbf{x}}   \newcommand{\vy}{\mathbf{y}}   \newcommand{\vz}{\mathbf{z}}
\newcommand{\vw}{\mathbf{w}}   \newcommand{\vu}{\mathbf{u}}   \newcommand{\vv}{\mathbf{v}}
\newcommand{\vb}{\mathbf{b}}   \newcommand{\vh}{\mathbf{h}}   \newcommand{\vr}{\mathbf{r}}
\newcommand{\vc}{\mathbf{c}}   \newcommand{\vf}{\mathbf{f}}   \newcommand{\vg}{\mathbf{g}}
\newcommand{\vp}{\mathbf{p}}   \newcommand{\vq}{\mathbf{q}}   \newcommand{\vs}{\mathbf{s}}
\newcommand{\va}{\mathbf{a}}   \newcommand{\vd}{\mathbf{d}}   \newcommand{\vt}{\mathbf{t}}
\newcommand{\vk}{\mathbf{k}}   \newcommand{\vm}{\mathbf{m}}   \newcommand{\vn}{\mathbf{n}}
\newcommand{\vzero}{\mathbf{0}} \newcommand{\vone}{\mathbf{1}}
% bold matrices
\newcommand{\vX}{\mathbf{X}}   \newcommand{\vY}{\mathbf{Y}}   \newcommand{\vZ}{\mathbf{Z}}
\newcommand{\vW}{\mathbf{W}}   \newcommand{\vH}{\mathbf{H}}   \newcommand{\vI}{\mathbf{I}}
\newcommand{\vA}{\mathbf{A}}   \newcommand{\vB}{\mathbf{B}}   \newcommand{\vC}{\mathbf{C}}
\newcommand{\vD}{\mathbf{D}}   \newcommand{\vP}{\mathbf{P}}   \newcommand{\vQ}{\mathbf{Q}}
\newcommand{\vU}{\mathbf{U}}   \newcommand{\vV}{\mathbf{V}}   \newcommand{\vS}{\mathbf{S}}
\newcommand{\vM}{\mathbf{M}}   \newcommand{\vK}{\mathbf{K}}   \newcommand{\vG}{\mathbf{G}}
% bold Greek
\newcommand{\vbeta}{\boldsymbol{\beta}}     \newcommand{\valpha}{\boldsymbol{\alpha}}
\newcommand{\vtheta}{\boldsymbol{\theta}}   \newcommand{\vmu}{\boldsymbol{\mu}}
\newcommand{\vSigma}{\boldsymbol{\Sigma}}   \newcommand{\vLambda}{\boldsymbol{\Lambda}}
\newcommand{\veps}{\boldsymbol{\varepsilon}} \newcommand{\vphi}{\boldsymbol{\phi}}
\newcommand{\vgamma}{\boldsymbol{\gamma}}   \newcommand{\vlambda}{\boldsymbol{\lambda}}
\newcommand{\vdelta}{\boldsymbol{\delta}}   \newcommand{\vnu}{\boldsymbol{\nu}}
\newcommand{\vxi}{\boldsymbol{\xi}}         \newcommand{\vpi}{\boldsymbol{\pi}}
\newcommand{\vomega}{\boldsymbol{\omega}}



% =============================================================================
\usepackage[margin=1in]{geometry}
\usepackage{helvet}
\usepackage[default]{lato}
\usepackage{mathpazo}
\usepackage[utf8]{inputenc}
\usepackage{amsmath,amssymb,amsthm}
\usepackage{mathtools}
\usepackage{graphicx}
\usepackage{booktabs}
\usepackage{multirow}
\usepackage{enumitem}
\usepackage{xcolor}
\usepackage{tcolorbox}
\tcbuselibrary{skins,breakable}
\usepackage{listings}
\usepackage[ruled,vlined]{algorithm2e}
\usepackage{hyperref}
\usepackage{fancyhdr}
\usepackage{titlesec}

\definecolor{blue}{RGB}{0,114,178}
\definecolor{red}{RGB}{213,94,0}
\definecolor{yellow}{RGB}{240,228,66}
\definecolor{green}{RGB}{0,158,115}
\hypersetup{colorlinks=true, linkcolor=blue, urlcolor=blue, citecolor=blue}

\titleformat{\section}{\Large\bfseries\color{blue}}{\thesection}{0.6em}{}
\titleformat{\subsection}{\large\bfseries}{\thesubsection}{0.6em}{}
\pagestyle{fancy}\fancyhf{}
\rhead{\footnotesize ENEI -- INEI \,\textperiodcentered\, PEU-CD 2026}
\lfoot{\footnotesize Machine Learning I \& II \,\textperiodcentered\, Alexander Quispe}
\rfoot{\footnotesize\thepage}
\renewcommand{\headrulewidth}{0.3pt}

\newtcolorbox{derivation}[1]{enhanced, breakable, colback=blue!3, colframe=blue, fonttitle=\bfseries,
  title={#1}, boxrule=0.6pt, left=6pt, right=6pt, top=4pt, bottom=4pt, arc=2pt}
\newtcolorbox{claim}[1]{enhanced, breakable, colback=white, colframe=black, fonttitle=\bfseries,
  title={#1}, boxrule=0.6pt, left=6pt, right=6pt, top=4pt, bottom=4pt, arc=2pt}
\newtcolorbox{keypoint}{enhanced, breakable, colback=green!6, colframe=green, boxrule=0.6pt,
  left=6pt, right=6pt, top=4pt, bottom=4pt, arc=2pt}
\newtcolorbox{caution}{enhanced, breakable, colback=red!5, colframe=red, boxrule=0.6pt,
  left=6pt, right=6pt, top=4pt, bottom=4pt, arc=2pt}
\newtcolorbox{task}[1]{enhanced, breakable, colback=yellow!12, colframe=yellow!60!black, fonttitle=\bfseries,
  title={#1}, boxrule=0.6pt, left=6pt, right=6pt, top=4pt, bottom=4pt, arc=2pt}

\lstset{
  basicstyle=\ttfamily\small, keywordstyle=\color{blue}\bfseries,
  commentstyle=\color{green}, stringstyle=\color{red},
  showstringspaces=false, breaklines=true, frame=single, rulecolor=\color{black!30},
  numbers=none, xleftmargin=6pt, framexleftmargin=6pt, language=Python
}
\newtheorem{theorem}{Theorem}
\newtheorem{lemma}{Lemma}
\newtheorem{proposition}{Proposition}
\theoremstyle{definition}
\newtheorem{definition}{Definition}
\newtheorem{exercise}{Exercise}

% =============================================================================
%  Notation — one convention for all lectures
%
%  scalars      x, y, w           plain italic
%  vectors      \vx \vy \vw       bold lower-case
%  matrices     \vX \vW \vSigma   bold upper-case
%  parameters   \vbeta \vtheta    bold Greek
%  transpose    \vX\T             \top
% =============================================================================
\newcommand{\R}{\mathbb{R}}
\newcommand{\E}{\mathbb{E}}
\newcommand{\Prob}{\mathbb{P}}
\DeclareMathOperator{\Var}{Var}
\DeclareMathOperator{\Cov}{Cov}
\DeclareMathOperator{\tr}{tr}
\DeclareMathOperator{\diag}{diag}
\DeclareMathOperator{\rank}{rank}
\DeclareMathOperator{\sign}{sign}
\DeclareMathOperator{\col}{col}
\DeclareMathOperator{\RSS}{RSS}
\DeclareMathOperator*{\argmin}{arg\,min}
\DeclareMathOperator*{\argmax}{arg\,max}
\newcommand{\norm}[1]{\left\lVert #1 \right\rVert}
\newcommand{\abs}[1]{\left\lvert #1 \right\rvert}
\newcommand{\inner}[2]{\left\langle #1,\, #2 \right\rangle}
\newcommand{\ind}[1]{\mathbf{1}\!\left\{#1\right\}}
\newcommand{\T}{^{\top}}
\newcommand{\inv}{^{-1}}
\newcommand{\half}{\tfrac{1}{2}}
\newcommand{\eps}{\varepsilon}
\newcommand{\Dset}{\mathcal{D}}
\newcommand{\Hset}{\mathcal{H}}
\newcommand{\Lcal}{\mathcal{L}}
\newcommand{\Ncal}{\mathcal{N}}

% bold vectors
\newcommand{\vx}{\mathbf{x}}   \newcommand{\vy}{\mathbf{y}}   \newcommand{\vz}{\mathbf{z}}
\newcommand{\vw}{\mathbf{w}}   \newcommand{\vu}{\mathbf{u}}   \newcommand{\vv}{\mathbf{v}}
\newcommand{\vb}{\mathbf{b}}   \newcommand{\vh}{\mathbf{h}}   \newcommand{\vr}{\mathbf{r}}
\newcommand{\vc}{\mathbf{c}}   \newcommand{\vf}{\mathbf{f}}   \newcommand{\vg}{\mathbf{g}}
\newcommand{\vp}{\mathbf{p}}   \newcommand{\vq}{\mathbf{q}}   \newcommand{\vs}{\mathbf{s}}
\newcommand{\va}{\mathbf{a}}   \newcommand{\vd}{\mathbf{d}}   \newcommand{\vt}{\mathbf{t}}
\newcommand{\vk}{\mathbf{k}}   \newcommand{\vm}{\mathbf{m}}   \newcommand{\vn}{\mathbf{n}}
\newcommand{\vzero}{\mathbf{0}} \newcommand{\vone}{\mathbf{1}}
% bold matrices
\newcommand{\vX}{\mathbf{X}}   \newcommand{\vY}{\mathbf{Y}}   \newcommand{\vZ}{\mathbf{Z}}
\newcommand{\vW}{\mathbf{W}}   \newcommand{\vH}{\mathbf{H}}   \newcommand{\vI}{\mathbf{I}}
\newcommand{\vA}{\mathbf{A}}   \newcommand{\vB}{\mathbf{B}}   \newcommand{\vC}{\mathbf{C}}
\newcommand{\vD}{\mathbf{D}}   \newcommand{\vP}{\mathbf{P}}   \newcommand{\vQ}{\mathbf{Q}}
\newcommand{\vU}{\mathbf{U}}   \newcommand{\vV}{\mathbf{V}}   \newcommand{\vS}{\mathbf{S}}
\newcommand{\vM}{\mathbf{M}}   \newcommand{\vK}{\mathbf{K}}   \newcommand{\vG}{\mathbf{G}}
% bold Greek
\newcommand{\vbeta}{\boldsymbol{\beta}}     \newcommand{\valpha}{\boldsymbol{\alpha}}
\newcommand{\vtheta}{\boldsymbol{\theta}}   \newcommand{\vmu}{\boldsymbol{\mu}}
\newcommand{\vSigma}{\boldsymbol{\Sigma}}   \newcommand{\vLambda}{\boldsymbol{\Lambda}}
\newcommand{\veps}{\boldsymbol{\varepsilon}} \newcommand{\vphi}{\boldsymbol{\phi}}
\newcommand{\vgamma}{\boldsymbol{\gamma}}   \newcommand{\vlambda}{\boldsymbol{\lambda}}
\newcommand{\vdelta}{\boldsymbol{\delta}}   \newcommand{\vnu}{\boldsymbol{\nu}}
\newcommand{\vxi}{\boldsymbol{\xi}}         \newcommand{\vpi}{\boldsymbol{\pi}}
\newcommand{\vomega}{\boldsymbol{\omega}}




\title{\textcolor{blue}{Lab 2 --- Ridge, Lasso and Cross-Validation}\\[4pt]
\large Machine Learning I \,\textperiodcentered\, Tutorial}
\author{Rodrigo Grijalba (teaching assistant) \,\textperiodcentered\, Alexander Quispe (instructor)}
\date{Wednesday, September 16, 2026 \,\textperiodcentered\, 19:00--22:00}

\begin{document}
\maketitle

\begin{keypoint}
\textbf{Goal of the session.} Implement ridge and lasso from the formulas in Lecture 3 --- closed form,
SVD shrinkage, soft thresholding, coordinate descent --- check each against scikit-learn, then choose
$\lambda$ properly and see what regularization does to a logistic regression that has no MLE.
\end{keypoint}

\section*{Materials}
\begin{itemize}[nosep]
  \item \texttt{lab\_02\_student.ipynb} / \texttt{lab\_02\_solution.ipynb}.
  \item Sources: CS 155 Set 2 (loss functions, effects of regularization, lasso vs.\ ridge).
  \item Data: \texttt{load\_diabetes} (regression) and \texttt{load\_breast\_cancer} (classification), both bundled with scikit-learn.
\end{itemize}

\section*{A note on conventions}
Every library defines the penalty slightly differently. Keep this table next to you.
\begin{center}\small
\begin{tabular}{lll}
\toprule
 & Objective & Relation to Lecture 3 \\
\midrule
Lecture 3, ridge & $\norm{\vy-\vX\vbeta}^2+\lambda\norm\vbeta^2$ & --- \\
\texttt{sklearn.linear\_model.Ridge(alpha)} & $\norm{\vy-\vX\vw}^2+\alpha\norm\vw^2$ & $\alpha=\lambda$ \\
Lecture 3, lasso & $\frac12\norm{\vy-\vX\vbeta}^2+\lambda\norm\vbeta_1$ & --- \\
\texttt{sklearn.linear\_model.Lasso(alpha)} & $\frac{1}{2N}\norm{\vy-\vX\vw}^2+\alpha\norm\vw_1$ & $\alpha=\lambda/N$ \\
\texttt{LogisticRegression(C)} & $\frac12\norm\vw^2+C\sum_i\log\bigl(1+e^{-\tilde y_i\vx_i\T\vw}\bigr)$ & $C=1/\lambda$ \\
\bottomrule
\end{tabular}
\end{center}
Both scikit-learn regressors fit the intercept separately and unpenalized, as Lecture 3 prescribes.

\section{Ridge from the Formula (35 min)}

\begin{task}{Task 1 --- closed form}
Standardize the diabetes features and centre $\vy$. Implement
$\hat\vbeta_\lambda=(\vX\T\vX+\lambda\vI)\inv\vX\T\vy$ with \texttt{np.linalg.solve}. For $\lambda\in\{0.1,1,10,100\}$,
compare with \texttt{Ridge(alpha=$\lambda$, fit\_intercept=False)} on the same standardized data. Agreement to $10^{-8}$.
\end{task}

\begin{task}{Task 2 --- the SVD picture}
Compute the thin SVD $\vX=\vU\vD\vV\T$ (\texttt{np.linalg.svd(X, full\_matrices=False)}). Verify numerically
that for $\lambda=10$
\[
\hat\vy_\lambda=\sum_{j=1}^p\vu_j\,\frac{d_j^2}{d_j^2+\lambda}\,\vu_j\T\vy
\]
matches $\vX\hat\vbeta_\lambda$ from Task 1. Print the ten shrinkage factors $d_j^2/(d_j^2+\lambda)$. Which
direction is shrunk most, and what is its $d_j$? Then plot $\mathrm{df}(\lambda)=\sum_jd_j^2/(d_j^2+\lambda)$
for $\lambda\in[10^{-2},10^4]$ on a log axis. Confirm $\mathrm{df}(0)=10$.
\end{task}

\begin{task}{Task 3 --- the ridge path}
Plot all ten $\hat\beta_j(\lambda)$ against $\lambda$ on a log axis, $\lambda\in[10^{-2},10^4]$. Every curve
must go to zero and none may cross zero exactly (ridge never produces an exact zero from a nonzero OLS
coefficient --- explain why with the SVD formula).
\end{task}

\section{Lasso from the Formula (50 min)}

\begin{task}{Task 4 --- soft thresholding}
Write \texttt{soft(b, lam)} $=\sign(b)\,(\abs b-\lambda)_+$, vectorized. Check the Lecture 3 example:
\texttt{soft([3, 0.5, -2], 1)} must return \texttt{[2, 0, -1]}, and the ridge shrinkage $b/(1+\lambda)$
must return \texttt{[1.5, 0.25, -1]}.
\end{task}

\begin{task}{Task 5 --- coordinate descent}
Implement the lasso by coordinate descent for the Lecture 3 objective $\frac12\norm{\vy-\vX\vbeta}^2+\lambda\norm\vbeta_1$:
\begin{lstlisting}
def lasso_cd(X, y, lam, iters=20000, tol=1e-10):
    n, p = X.shape
    beta = np.zeros(p)
    r = y.copy()                          # full residual y - X beta
    for _ in range(iters):
        for j in range(p):
            r_j = r + X[:, j] * beta[j]   # partial residual, coordinate j removed
            b_new = soft(X[:, j] @ r_j, lam) / (X[:, j] @ X[:, j])
            r = r_j - X[:, j] * b_new     # update the residual incrementally
            beta[j] = b_new
    return beta
\end{lstlisting}
Add a stopping rule (sweep until no coordinate moves more than \texttt{tol}; correlated columns make plain coordinate descent slow), then compare with \texttt{Lasso(alpha=lam/N, fit\_intercept=False)} for
$\lambda\in\{50,200,800\}$ on standardized, centred data. Which coefficients are exactly zero? Verify the KKT
condition on the solution: $\abs{\vx^{(j)\top}\vr}=\lambda$ for every active $j$ and $\le\lambda$ for every inactive $j$.
\end{task}

\begin{task}{Task 6 --- the lasso path and $\lambda_{\max}$}
Compute $\lambda_{\max}=\norm{\vX\T\vy}_\infty$. Show that \texttt{lasso\_cd(X, y, 1.01 * lam\_max)} returns all
zeros and \texttt{lasso\_cd(X, y, 0.99 * lam\_max)} does not. Then plot the full path with
\texttt{sklearn.linear\_model.lasso\_path}. In what order do features enter? Compare with the ridge path of
Task 3 on the same axes.
\end{task}

\section{Choosing $\lambda$ (35 min)}

\begin{task}{Task 7 --- cross-validation with error bars}
For 40 values of $\lambda$ on a log grid, compute the 10-fold CV MSE of the lasso and its standard error
across folds (use \texttt{KFold(10, shuffle=True, random\_state=155)}; standardize \emph{inside} each fold
with a \texttt{Pipeline}). Plot the curve with error bars. Mark $\hat\lambda_{\min}$ and the one-standard-error
choice $\hat\lambda_{1\text{se}}$. How many nonzero coefficients does each have?
\end{task}

\begin{task}{Task 8 --- optimism, measured}
Lecture 3 proves $\E[\text{test}]-\E[\text{train}]=2\sigma^2\,\mathrm{df}/N$ for a linear smoother. Test it on ridge:
for $\lambda\in\{0.1,1,10,100,1000\}$, compute (i) training MSE, (ii) 10-fold CV MSE, (iii)
$2\hat\sigma^2\,\mathrm{df}(\lambda)/N$ with $\hat\sigma^2$ from the OLS residuals. Tabulate
(ii)$-$(i) against (iii). They should track each other in size and in trend.
\end{task}

\section{Regularized Logistic Regression (40 min)}

\begin{task}{Task 9 --- separable data has no MLE}
Generate two well-separated Gaussian clouds in $\R^2$ (30 points each, means $(-2,0)$ and $(2,0)$, unit
variance). Fit \texttt{LogisticRegression(penalty=None, max\_iter=10000)} and print $\norm{\hat\vw}$ and the
fitted probabilities. Then fit with \texttt{C} $\in\{10^3,10^2,10,1,0.1\}$ and plot $\norm{\hat\vw}$ against $C$
on a log axis. Explain the shape using Lecture 2, slide ``When the MLE Does Not Exist''.
\end{task}

\begin{task}{Task 10 --- a real classifier, evaluated properly}
On \texttt{load\_breast\_cancer} ($N=569$, $p=30$), with an 80/20 split:
\begin{enumerate}[nosep]
  \item Fit $\ell_2$-regularized logistic regression with $C$ chosen by 5-fold CV on the training set.
  \item Report accuracy, precision, recall, $F_1$ and the confusion matrix on the test set at threshold $0.5$.
  \item Plot the ROC curve and compute the AUC two ways: with \texttt{roc\_auc\_score}, and directly as the
        fraction of (positive, negative) test pairs with the positive scored higher. They must agree.
  \item Suppose a false negative (missed malignancy) costs 10 times a false positive. Compute
        $\tau^\star=c_{\text{FP}}/(c_{\text{FP}}+c_{\text{FN}})$, re-threshold, and report the new confusion matrix and expected cost.
        Compare with the cost at $\tau=0.5$.
\end{enumerate}
\end{task}

\section{Exercises (to submit before Lab 3)}
\begin{exercise}
Hoerl--Kennard. With $\vX\T\vX=\vV\vD^2\vV\T$ and $\vgamma=\vV\T\vbeta$, show that the total mean squared error
$\E\norm{\hat\vbeta_\lambda-\vbeta}^2$ of ridge equals
$\sum_j\bigl[\sigma^2d_j^2+\lambda^2\gamma_j^2\bigr]/(d_j^2+\lambda)^2$, differentiate at $\lambda=0$, and
conclude that its derivative is negative --- so some $\lambda>0$ beats OLS.
\end{exercise}
\begin{exercise}
Extend \texttt{lasso\_cd} to the elastic net $\frac12\norm{\vy-\vX\vbeta}^2+\lambda_1\norm\vbeta_1+\frac{\lambda_2}{2}\norm\vbeta^2$
by deriving the coordinate update (it is soft thresholding divided by $\norm{\vx^{(j)}}^2+\lambda_2$).
Check against \texttt{ElasticNet}.
\end{exercise}
\begin{exercise}
Repeat Task 7 with 20 different random fold assignments. Report the mean and standard deviation of
$\hat\lambda_{\min}$ and of $\hat\lambda_{1\text{se}}$ (on the log scale). Which choice is more stable, and why
would you expect that from the shape of the CV curve?
\end{exercise}

\section*{Submission}
One executed notebook, \texttt{lab02\_lastname\_firstname.ipynb}, to the course drive before
Wednesday, September 23, 18:00.

\end{document}
